Este trabalho investigou a execução de Busca em Largura em GPU usando CUDA e
NVSHMEM, tendo o Graph500 como referência metodológica e experimental. A partir
desse benchmark, foram desenvolvidas duas abordagens: uma implementação
CUDA/NVSHMEM compatível com o fluxo do Graph500 e uma versão refatorada em forma
de biblioteca de travessia de grafos executada majoritariamente em GPU.

Os resultados mostram que a geração R-MAT se beneficia fortemente do paralelismo
da GPU nas maiores escalas avaliadas. No computador pessoal, em
\texttt{SCALE=22}, a geração em GPU foi aproximadamente $7{,}4$ vezes mais
rápida que a geração em CPU. No ambiente GB200, em \texttt{SCALE=29}, a biblioteca
GPU reduziu o tempo de geração de 31,13 s para 2,738 s em relação à referência em
CPU.

Para a BFS, a versão CUDA/NVSHMEM híbrida superou a referência em CPU apenas nas
maiores escalas do computador pessoal, indicando que o custo fixo de coordenação
e sincronização ainda é significativo em entradas menores. A biblioteca GPU
apresentou ganhos mais expressivos, especialmente por manter o laço principal da
BFS no dispositivo. No computador pessoal, em \texttt{SCALE=21}, a biblioteca
alcançou 2,138 GTEPS, contra 0,236604 GTEPS da versão híbrida. No GB200, em
\texttt{SCALE=29}, a biblioteca alcançou 6,180 GTEPS, contra 1,570 GTEPS da
referência em CPU.

Esses resultados reforçam a hipótese de que, em travessias distribuídas de
grafos, mover apenas a expansão da fronteira para a GPU não é suficiente. A
redução da participação recorrente da CPU na coordenação entre níveis é um fator
importante para melhorar o desempenho, sobretudo conforme a escala do grafo e o
número de elementos de processamento aumentam.

Como limitações, a BFS permanece sensível à irregularidade do grafo, à contenção
em operações atômicas, à comunicação remota e ao uso de memória. Além disso, a
biblioteca ainda representa uma etapa de refatoração em andamento. Trabalhos
futuros incluem aprimorar a escalabilidade entre múltiplas GPUs, otimizar a
comunicação entre PEs, reduzir contenção em atualizações remotas, ampliar a
validação experimental e consolidar uma interface de biblioteca mais estável para
outros algoritmos de grafos.
